\chapter{The Risk Grid}\label{ch:pl:the-risk-grid}

The nightly risk batch --- every trade, every historical scenario --- starts at 04:30 and must be done by 06:00, when the morning's risk
reports are built. It finished at 06:51. Doubling the machines moved the finish to 06:25 and no earlier, because the last hour belonged to
one task waiting for one trade that priced two hundred times slower than its estimate, and the scheduler, which believed the estimate,
had started that task late. Nothing was broken; the grid was planned by trade counts, not by cost. This chapter plans, runs and repairs
such a grid on the simulated cluster of chapter~13, with the costs of the real pricing library measured on this laptop.

\section{Nightly revaluation at scale}

\begin{definition}[Risk grid, batch window]\label{def:pl:the-risk-grid:grid}
A \emph{risk grid}\index{risk grid} is the infrastructure that revalues a firm's trades under many market scenarios on many machines: the
plan that cuts the work into tasks, the scheduler and cluster that run them, and the store that collects the results for the risk engine to
aggregate. Its \emph{batch window}\index{batch window} is the interval between the moment its inputs are ready (the end-of-day positions and
market snapshot) and the moment its outputs are needed.
\end{definition}

Book~6's risk engine (chapter~29) computes a book's P\&L vectors under historical scenarios by full revaluation, a revaluation grid or a
delta-gamma approximation, and aggregates them into VaR and expected shortfall along the risk hierarchy. At the size of a firm the
computation is a grid of trades by scenarios --- here 105\,000 trades and 1\,000 scenarios, 105~million pricing calls --- and the engine
becomes a scheduling problem.

The chapter's book: 60\,000 swaps and 20\,000 swaptions (Book~6's rates instruments), 20\,000 equity options and 5\,000 autocallables
priced by Monte Carlo with the library's 100\,000 paths (Book~5). One pricing call of each, measured on this laptop with the real library,
costs 14.5, 10.2 and 5.1~microseconds for the first three and 22.2~milliseconds for an autocallable: the 5\,000 autocallables are 99\% of
the book's 112.2~seconds of pricing per scenario. A task costs a fixed five seconds (loading the snapshot and the trades, writing the results)
plus its pricing. The grid has 32 cores (eight nodes of four), and the window is 90~minutes.

\section{Fanning out scenarios}

\begin{definition}[Scenario fan-out, task granularity]\label{def:pl:the-risk-grid:fanout}
\emph{Scenario fan-out}\index{scenario fan-out} is the splitting of the scenarios into batches so that one trade batch's revaluation runs as
several tasks in parallel. \emph{Task granularity}\index{task granularity} is the size of the resulting tasks --- here the number of scenarios
in a batch, the trades being grouped into batches of about half a second of pricing per scenario --- which trades scheduling freedom against
the fixed cost of each task.
\end{definition}

\begin{omfigure}
\begin{tikzpicture}[font=\scriptsize,>=Stealth]
\foreach \r in {0,...,4} {\foreach \c in {0,...,7} {\draw[fill=omDef!12] (\c*1.1,-\r*0.55) rectangle ++(1.0,0.45);}}
\foreach \c in {0,...,7} {\node at (\c*1.1+0.5,0.75) {$s_{\c}$};}
\foreach \r/\l in {0/{swaps 1},1/{swaps 2},2/{\dots},3/{autocallables},4/{slow trade}} {\node[anchor=east] at (-0.1,-\r*0.55+0.22) {\l};}
\foreach \c in {0,...,15} {\draw[fill=omThm!45] (\c*0.55,-4*0.55) rectangle ++(0.5,0.45);}
\node[anchor=west,align=left] at (9.0,-1.0) {a task = one trade batch\\ $\times$ one scenario batch};
\node[anchor=west,align=left,omThm] at (9.0,-2.0) {the slow trade: isolated,\\ its scenarios fanned out finely};
\node at (4.3,1.2) {scenario batches};
\end{tikzpicture}

\omcaption{The grid as tasks: rows are trade batches grouped by estimated cost, columns scenario batches; every cell is a task. A cost-aware
plan gives a trade measured slow its own row, with finer columns, so that no single task holds the window.}\label{fig:pl:the-risk-grid:tasks}
\end{omfigure}

\omcode{code/firm/riskgrid/firm_riskgrid.py}{77}{92}{Planning the grid: trade batches by estimated cost, scenario batches of a chosen size,
and trades too costly for one task isolated and fanned out more finely.}\label{lst:pl:the-risk-grid:plan}

With 1\,000 scenarios in a batch there are 220 tasks, each a trade batch across all scenarios: few, long, and impossible to balance. With
five scenarios there are 44\,000, and their fixed costs --- 44\,000 times five seconds --- are larger than the pricing itself.
\Cref{fig:pl:the-risk-grid:granularity} is the whole curve: on 32 cores the \omterm{def:pl:distributed-compute-and-schedulers:makespan}{makespan} falls from 2.35~hours at 1\,000 scenarios a task to
1.19 at 100, and rises again to 2.93 at five; only batches of 20 to 250 scenarios meet the window, 20 by sixteen seconds.

\begin{omfigure}
\begin{tikzpicture}
\begin{semilogxaxis}[width=12cm,height=6cm,xlabel={scenarios per task (log scale)},ylabel={makespan (hours)},xmin=4,xmax=1200,ymin=0,ymax=3.6,
  grid=major,grid style={gray!20},tick label style={font=\small},label style={font=\small},table/col sep=comma,
  legend style={font=\footnotesize,at={(0.5,1.03)},anchor=south,legend columns=2}]
\addplot[color=omThm,very thick,mark=*] table[x=scen_batch,y=naive_h]{figdata/platforms/20-the-risk-grid/sweep_32.csv};
\addlegendentry{naive, 32 cores}
\addplot[color=omDef,very thick,mark=square*] table[x=scen_batch,y=aware_h]{figdata/platforms/20-the-risk-grid/sweep_32.csv};
\addlegendentry{cost-aware, 32 cores}
\addplot[color=omThm,thick,dashed,mark=o] table[x=scen_batch,y=naive_h]{figdata/platforms/20-the-risk-grid/sweep_64.csv};
\addlegendentry{naive, 64 cores}
\addplot[color=omDef,thick,dashed,mark=square] table[x=scen_batch,y=aware_h]{figdata/platforms/20-the-risk-grid/sweep_64.csv};
\addlegendentry{cost-aware, 64 cores}
\draw[black,thick,dotted] (axis cs:4,1.5) -- (axis cs:1200,1.5);
\end{semilogxaxis}
\end{tikzpicture}

\omcaption{\omterm{def:pl:distributed-compute-and-schedulers:makespan}{Makespan} of the nightly grid on 32 cores (solid) and 64 (dashed) against the number of scenarios per task, for a plan that trusts the cost estimates
(naive) and one that isolates the trade measured slow last night (cost-aware); the dotted line is the 90-minute window. Coarse tasks cannot be balanced; fine tasks drown in their
fixed costs. Data: \texttt{fig\_riskgrid.py}, with costs from \texttt{bench\_riskgrid.py}.}\label{fig:pl:the-risk-grid:granularity}
\end{omfigure}

\section{Task granularity, failure and retry}

The slow trade makes coarse tasks worse still. One autocallable's engine falls back to daily steps and prices at 4.88~seconds, 220 times its
estimate. In the naive plan its trade batch across all 1\,000 scenarios is a task of more than 5\,300~seconds that the scheduler, sorting by
estimate, starts like any other: the grid finishes 8\,468~seconds after 04:30, at 06:51; on 64 cores, at 06:25 (6\,922~seconds). A cost-aware
plan uses last night's measured cost: the slow trade gets tasks of its own, 122 scenarios each, and the 1\,000-scenario plan finishes in
4\,123~seconds on 32 cores (05:38); the best granularity, 250 scenarios, in 3\,846 (05:34).

\omcode{code/firm/riskgrid/firm_riskgrid.py}{105}{122}{Running a plan: failures handled per task (retried, then lost) or per trade
(quarantined), then the tasks scheduled longest estimate first on the simulated cluster.}\label{lst:pl:the-risk-grid:run}

Failures need the same thinking. One swaption's volatility is missing, and it fails every time. If a failing trade fails its task, the
scheduler retries the task twice and then gives up: the grid loses every result of every task that held the trade --- 38.3~million of the
105~million cells, more than a third of the book --- and burns 1\,100 core-seconds on retries. If the pricing loop catches the failure per trade, as
Book~6's engine does (a failing trade gives a missing value and an entry in an error report, never a stopped run), the grid loses the
failing trade's 1\,000 cells, quarantines it with a report, and finishes 60~seconds sooner.

\begin{table}[ht]
\centering\small
\begin{tabular}{lrrr}
\toprule
failure handling & \omterm{def:pl:distributed-compute-and-schedulers:makespan}{makespan} (s) & core-seconds wasted & results lost \\
\midrule
per task (two retries, then lost) & 4\,095 & 1\,100 & 38\,263\,000 \\
per trade (quarantined with a report) & 4\,035 & 0 & 1\,000 \\
\bottomrule
\end{tabular}
\caption{One trade that always fails, in the cost-aware plan with 100 scenarios per task on 32 cores.}\label{tab:pl:the-risk-grid:failures}
\end{table}

\section{Adjoint sensitivities in production}

The nightly batch also computes sensitivities. Bumping a swap's eight curve pillars up and down costs sixteen repricings; the adjoint of
Book~4's tape (\texttt{firm.aad}, chapter~28 there) records one valuation and sweeps back once, and gives all eight derivatives --- equal to the
bumped ones to $10^{-9}$ in the tests (\cref{lst:pl:the-risk-grid:adjoint}). On this laptop the bumped deltas of a swap take 388~microseconds
and the adjoint 66, 5.9 times faster, with a Python tape; a compiled tape keeps the adjoint's cost at a few valuations whatever the number
of pillars. Over the book's 60\,000 swaps, bumping needs 1.02~million pricing calls, the adjoint 60\,000 recorded valuations.

\omcode{code/firm/riskgrid/firm_riskgrid.py}{126}{141}{A swap's value on a pillar curve written once on the tape; one reverse sweep gives
its derivative to every pillar.}\label{lst:pl:the-risk-grid:adjoint}

\section{Intraday risk and incremental revaluation}

\begin{definition}[Incremental revaluation]\label{def:pl:the-risk-grid:incremental}
\emph{Incremental revaluation}\index{incremental revaluation} recomputes only the results whose inputs changed: each result is stored with the
version of the trade and of the market data it was computed from, and a rerun reprices only the trades whose key is not in the store.
\end{definition}

Intraday, the scenarios and the base snapshot are the morning's, and most trades have not changed. With results keyed by (trade, trade
version, market-data version), an intraday rerun after 1\,500 trades are amended or booked reprices those 1\,500 under the 1\,000 scenarios:
1.5~million pricing calls instead of 105~million. A new market snapshot changes every key; for that the real-time system of chapter~18 carries
sensitivities forward, and the grid waits for the night.

\begin{omfigure}
\begin{tikzpicture}
\begin{axis}[width=11.5cm,height=5.4cm,xbar,bar width=9pt,xmode=log,xmin=1e4,xmax=5e9,log origin=infty,
  symbolic y coords={intraday incremental,intraday full,swap deltas adjoint,swap deltas bumped,nightly grid},ytick=data,
  enlarge y limits=0.15,xlabel={pricing calls (log scale)},tick label style={font=\small},label style={font=\small},
  grid=major,grid style={gray!20},nodes near coords,nodes near coords style={font=\scriptsize,anchor=west},point meta=explicit symbolic]
\addplot[fill=omDef!45,draw=omDef] coordinates {(105105000,nightly grid)[105.1 million] (1020000,swap deltas bumped)[1.02 million]
  (60000,swap deltas adjoint)[60\,000 valuations] (105000000,intraday full)[105 million] (1500000,intraday incremental)[1.5 million]};
\end{axis}
\end{tikzpicture}

\omcaption{Pricing calls of the nightly grid, of the swaps' pillar deltas by bumping and by the adjoint, and of an intraday rerun done in full
and incrementally after 1\,500 trades changed.}\label{fig:pl:the-risk-grid:calls}
\end{omfigure}

\begin{example}[06:51 against 06:00]\label{ex:pl:the-risk-grid:batch}
Planned by trade batches across all scenarios and trusting its estimates, the grid finishes at 06:51 on 32 cores and at 06:25 on 64: more
cores shorten the queue, not the one long task. Cutting the scenarios into batches of 100 brings the naive plan to 05:41; isolating the slow
trade by its measured cost brings every granularity from 20 to 1\,000 scenarios inside the window, the best at 05:34 (250 scenarios a task).
Quarantining the failing trade instead of failing its tasks keeps more than a third of the book's results. The pricing calls saved elsewhere are larger
still: the adjoint replaces 1.02~million bumped repricings of swaps by 60\,000 recorded valuations, and an incremental intraday rerun makes
1.5~million calls where a full one makes 105~million.
\end{example}

\section{Tutorial: planning a nightly grid}\label{tut:pl:the-risk-grid:tut}

\begin{tutorial}
\textbf{Goal.} Measure pricing costs, plan the grid at several granularities, repair it for a slow and a failing trade, and count the calls
adjoints and incremental runs save. \textbf{End state:} \cref{fig:pl:the-risk-grid:granularity} and \cref{tab:pl:the-risk-grid:failures}.
\begin{tutsteps}
\item \textbf{Costs}: \texttt{bench\_riskgrid.py}: one pricing call per kind, bumped and adjoint swap deltas.
\item \textbf{Plan and run}: \texttt{pl\_riskgrid.grid(scen\_batch, cost\_aware)}; \texttt{sweep()}.
\item \textbf{Failures}: \texttt{failure\_policies()}.
\item \textbf{Adjoints}: \texttt{firm\_riskgrid.swap\_pv\_adjoint} against the library's bumped sensitivities.
\item \textbf{Intraday}: \texttt{intraday()}.
\end{tutsteps}
\textbf{What to change next.} Give the scheduler the measured costs of every trade and see how close to the lower bound the grid comes; add
the sensitivities' repricings to the grid's work and find the new best granularity.
\end{tutorial}

\section{Build: the risk grid}\label{bld:pl:the-risk-grid:riskgrid}

\begin{build}
\bfield{Purpose} The nightly revaluation planned by cost, run to its window, robust to slow and failing trades, and cheaper by adjoints and
incremental reruns.
\bfield{Interface} \texttt{GridTrade}, \texttt{plan}, \texttt{run -> GridRun}, \texttt{swap\_pv\_adjoint}, \texttt{ResultCache},
\texttt{cells}, \texttt{lower\_bound}, \texttt{hours}.
\bfield{Rules} Plans use costs measured on the previous run where they exist; a trade's failure never fails its task; results carry their trade
and market-data versions; adjoints are checked against bumps.
\bfield{Acceptance tests} \texttt{code/firm/riskgrid/tests/}: batching and isolation of a heavy trade; the two failure policies and the lower
bound; adjoint pillar deltas equal to the library's bumped sensitivities; the result cache's keys.
\bfield{Stretch} Results written to chapter~11's \omterm{def:pl:backtest-engine-architecture:results}{results store}; the plan fed by a cost model learned from past runs; checkpointed adjoints
through Monte Carlo paths (Book~4, chapter~28).
\end{build}

\begin{omsources}
\item One Quant Book~4, chapter~28 (adjoints); One Quant Book~5, chapter~28, and Book~6, chapter~29 (the library, the risk engine); chapter~13 of this book (scheduling).
\end{omsources}

\section{Exercises}

\begin{exercise}[$\star$]\label{exo:pl:the-risk-grid:1}
How many tasks does the plan make with 100 scenarios per task, and how much fixed cost do they add?
\end{exercise}

\begin{exercise}[$\star$]\label{exo:pl:the-risk-grid:2}
Why does doubling the cores move the naive plan's finish from 06:51 only to 06:25?
\end{exercise}

\begin{exercise}[$\star$]\label{exo:pl:the-risk-grid:3}
How many pricing calls do bumped pillar deltas of a swap need, and how many valuations does the adjoint need?
\end{exercise}

\begin{exercise}[$\star\star$]\label{exo:pl:the-risk-grid:4}
Why does the \omterm{def:pl:distributed-compute-and-schedulers:makespan}{makespan} rise again below 50 scenarios per task?
\end{exercise}

\begin{exercise}[$\star\star$]\label{exo:pl:the-risk-grid:5}
Why does failing a task for one trade lose 38 million results?
\end{exercise}

\begin{exercise}[$\star\star$]\label{exo:pl:the-risk-grid:6}
When does \omterm{def:pl:the-risk-grid:incremental}{incremental revaluation} save nothing, and what covers that case?
\end{exercise}

\begin{exercise}[$\star\star\star$]\label{exo:pl:the-risk-grid:7}
\emph{Coding.} Run the sweep on 64 cores. Which granularities meet the window with each plan, and which is best?
\end{exercise}

\begin{exercise}[$\star\star\star$]\label{exo:pl:the-risk-grid:8}
\emph{Find the flaw.} ``The batch was late, so we bought twice the machines.''
\end{exercise}

\section{Problem: 06:51 Against 06:00}

\begin{problem}[{Weekend problem --- a grid that misses its window}]\label{pb:pl:the-risk-grid:1}
The chapter's book, costs, cluster and plans.

\medskip\noindent\textbf{Part I --- The grid.}
\begin{enumerate}
\item What is the grid's work, in trades, scenarios and pricing calls?
\item Which trades dominate its cost, and why?
\item What does a task cost besides its pricing?
\item What is the window, and on how many cores?
\item What is the lower bound on 32 cores?
\end{enumerate}

\medskip\noindent\textbf{Part II --- Granularity.}
\begin{enumerate}[resume]
\item What is the \omterm{def:pl:distributed-compute-and-schedulers:makespan}{makespan} at 1\,000, 100 and 5 scenarios a task?
\item Which granularities meet the window with the naive plan?
\item What does the slow trade do to coarse tasks, and why does the scheduler not help?
\item What does the cost-aware plan change?
\item What is the best cost-aware finish?
\end{enumerate}

\medskip\noindent\textbf{Part III --- Failures and savings.}
\begin{enumerate}[resume]
\item What happens to a failing trade per task and per trade?
\item How many results and core-seconds does each policy lose?
\item How fast are adjoint swap deltas against bumped ones, here?
\item How many pricing calls does the adjoint save over the book's swaps?
\item What does an incremental intraday rerun cost?
\end{enumerate}

\medskip\noindent\textbf{Part IV --- The verdict.}
\begin{enumerate}[resume]
\item State the \emph{named result}: the batch's \omterm{def:pl:distributed-compute-and-schedulers:makespan}{makespan} as a function of \omterm{def:pl:the-risk-grid:fanout}{task granularity} and cluster size, the granularity that meets the
window, and the pricing calls saved by adjoint sensitivities and by incremental intraday revaluation.
\item What would you change first in the grid's planning?
\item What would you monitor every night?
\item When would you buy more machines?
\item In one sentence: what decides when a nightly batch ends?
\end{enumerate}
\end{problem}

\section{Interview questions}

\begin{interviewq}[$\star$ \iqroles{developer}]\label{iq:pl:the-risk-grid:1}
Your nightly risk batch misses its window. What do you look at first?
\end{interviewq}

\begin{interviewq}[$\star\star$ \iqroles{developer}]\label{iq:pl:the-risk-grid:2}
How would you split a million trades by a thousand scenarios into tasks?
\end{interviewq}

\begin{interviewq}[$\star\star$ \iqroles{developer, researcher}]\label{iq:pl:the-risk-grid:3}
Why use adjoint differentiation for risk, and where does it not help?
\end{interviewq}

\begin{interviewq}[$\star\star$ \iqroles{developer}]\label{iq:pl:the-risk-grid:4}
One trade fails to price every night. What should the grid do with it?
\end{interviewq}

\begin{interviewq}[$\star\star\star$ \iqroles{developer}]\label{iq:pl:the-risk-grid:5}
Design the nightly and intraday revaluation infrastructure of a bank's trading book.
\end{interviewq}
